\chapter*{국문요약}
\addcontentsline{toc}{chapter}{국문요약}
\begin{center}\large\bfseries\thesistitlekr\end{center}

표면 마무리 공정에서 부품 전체의 균일 처리와 지정된 영역 주변의 국소 재작업은 서로 다른 요구를 가진다. 국소 폴리싱에서는 대상 영역에 맞추어 공구 이동, 접촉 힘, 작업 시간을 함께 구성해야 한다. 본 연구는 현재 장비로 미세 스크래치를 자율 판별하는 대신 사람이 처리 영역을 지정하는 설정을 사용한다. 공간 지시 이후에도 남는 동작·힘·시간 구성 문제에 대해, 규칙 기반 전략과 교시 재생에 비해 관측 조건부 학습이 어떤 수행 가치를 추가하는지 검토한다.

추적 기반 교시 장치로 공구 위치·자세, 3축 힘, RGB 영상을 시간적으로 대응하여 기록한다. 얼룩 중심 상대위치를 학습하고 얼룩 방향을 학습 기준에 정렬한다. Flow Matching 정책은 영상, 로봇 상태, 측정 힘 이력으로 위치·자세·목표 힘 시퀀스를 생성한다. 공통 실행 계층은 좌표 복원, 계획 연결, 명령 후처리를 수행하고 기존 하위 힘제어기가 접촉을 조절한다. 스킬 생성, 접촉 제어, 절차 완료 판정의 역할을 분리한다.

E1/E2는 방향 정규화 없이 평가한다. E1은 기준조건에서 규칙 기반 경로, 힘을 포함한 교시 재생, 힘 조건부 정책을 비교한다. 모든 실행에서 요구된 부하 편도 통과가 관측되었다. 기존 완료율 차이는 상승 이탈의 높이 기준에 민감하였고, 학습 정책의 초과 힘 적분은 다른 방법보다 컸다. 동일 기록으로 판정 민감도를 분석하고 실제 시편 높이 변화 시험은 계획으로 구분하였다. E2는 목표 힘 학습과 측정 힘 관측을 분리한다. 힘 관측 정책의 작업자 종료 구간이 더 짧았으나, 동작만 학습한 조건의 한 시행은 접촉 해제로 전환하지 못해 작업자가 중단하였다. 결과는 동작--힘 표현의 실행 가능성과 단계 진행의 차이를 보여주며, 정책의 보편적 우위나 표면 품질 향상을 입증하지는 않는다. 본 연구는 국소 작업에서 규칙·재생·학습의 활용 범위를 검토할 수 있는 구현과 평가 근거를 제공한다.

\vspace{0.4cm}
\noindent\textbf{주요어:} 국소 로봇 폴리싱, 인간 작업 지정, 동작--힘 스킬, 모방학습, Flow Matching, 접촉 작업 평가
